% Figure: the convex function class and its defining chord property.
% A convex curve with a chord lying above the graph between two points;
% the tangent line lying below the graph; the midpoint gap that is the
% content of Jensen's inequality.
\begin{figure}[htbp]
\centering
\resizebox{0.66\linewidth}{!}{%
\begin{tikzpicture}[scale=1.0,
  ax/.style={->, draw=engblack, thick},
  curve/.style={draw=engblue, very thick},
  chord/.style={draw=engred, thick},
  tang/.style={draw=engorange, thick, dash pattern=on 4pt off 2pt},
  pt/.style={circle, fill=engblack, inner sep=1.2pt},
  ann/.style={font=\footnotesize, align=center},
]
  % Convex curve f(x) = 0.45 x^2 on x in [-0.4, 3.4]; display scale 1:1.
  \draw[ax] (-0.4, 0) -- (4.0, 0) node[below right, ann] {$x$};
  \draw[ax] (0, -0.3) -- (0, 5.0) node[left, ann] {$f(x)$};
  \draw[curve, domain=-0.3:3.3, samples=60, smooth]
    plot (\x, {0.45*\x*\x});
  % Two points a = 0.5, b = 3.0.
  \def\ax{0.5}\def\bx{3.0}
  \pgfmathsetmacro{\ay}{0.45*\ax*\ax}
  \pgfmathsetmacro{\by}{0.45*\bx*\bx}
  % Chord from (a, f(a)) to (b, f(b)) lies above the graph.
  \draw[chord] (\ax, {\ay}) -- (\bx, {\by});
  \node[pt] at (\ax, {\ay}) {};
  \node[pt] at (\bx, {\by}) {};
  \node[ann, anchor=north east, text=engred] at ({(\ax+\bx)/2 + 0.5}, {(\ay+\by)/2 + 0.55})
    {chord lies\\above graph};
  % Midpoint of the chord vs the function at the midpoint (Jensen gap).
  \pgfmathsetmacro{\mx}{(\ax+\bx)/2}
  \pgfmathsetmacro{\mchord}{(\ay+\by)/2}
  \pgfmathsetmacro{\mfun}{0.45*\mx*\mx}
  \draw[engblack!55, thin, dotted] (\mx, {\mfun}) -- (\mx, {\mchord});
  \node[ann, anchor=west, fill=white, inner sep=1pt]
    at ({\mx + 0.08}, {(\mfun+\mchord)/2})
    {Jensen gap\\$\tfrac{f(a)+f(b)}{2}\!-\!f\!\left(\tfrac{a+b}{2}\right)$};
  % Tangent at x = 1.6 lies below the graph.
  \def\tx{1.6}
  \pgfmathsetmacro{\ty}{0.45*\tx*\tx}
  \pgfmathsetmacro{\tslope}{0.9*\tx}
  \draw[tang, domain=0.4:3.0, samples=2]
    plot (\x, {\ty + \tslope*(\x - \tx)});
  \node[ann, anchor=north west, text=engorange] at (2.0, 0.7)
    {tangent lies\\below graph};
\end{tikzpicture}%
}%
\caption[The convex function class and its chord property]{A convex
function has two equivalent defining properties drawn here. The chord
joining any two points on the graph (red) lies on or above the graph
between them. The tangent at any point (orange dashed) lies on or below
the graph. The gap between the chord midpoint and the function at the
midpoint is the content of Jensen's inequality, that the average of the
function exceeds the function of the average for a convex $f$. A convex
function on an interval has at most one local minimum, which is also
its global minimum, and that single guarantee is why convexity is the
property optimisation algorithms most want their objective to
possess.}
\label{fig:vol02:ch02:convex-classes}
\end{figure}
